\chapter{The archive format}
\label{cap:formato}

\section{Files and messages}
\label{sec:file}

The archive consists of one GRIB2 file per validity time. Every hourly file contains 90 variables
in 246 GRIB messages, always the same ones. The list comes from a single registry in the
converter, and every file produced is checked against that list, so a field that could not be
computed does not go unnoticed. File names follow the pattern
\begin{center}
\campo{ailam-an-cima-3km-{year}-{year}-1h-v1-{YYYYMMDD}{HH}.grib},
\end{center}
\noindent in which the year appears twice with the same value and the final part is the validity
time.

The messages use grid definition template 3.10 (Mercator), originating centre \campo{ecmf} and
\campo{grid_ccsds} packing at 24 bits per value. Every parameter carries the ECMWF paramId and
shortName \parencite{paramdb}. There are 89 distinct shortNames for 90 variables, because
\campo{z} appears twice, as geopotential on the pressure levels and as surface geopotential,
i.e.\ the orography. Fields defined only over part of the domain carry a GRIB bitmap, and points
outside the domain of definition are missing points, not zeros (\cref{sec:mascheratura}).

An hourly file is \SI{737}{MB} on average over the 184 files of September 2024, and
\SI{661}{MB} over the 8 files at 00 UTC among them; over the 476 files of February and March 2024
the mean is \SI{765}{MB}. This gives between \SI{17.7}{GB} and \SI{18.4}{GB} per day and between
\SI{6.5}{TB} and \SI{6.7}{TB} per year.

\section{Geometry}
\label{sec:geometria}

The grid declares the shape of the Earth with \campo{shapeOfTheEarth=1} and radius
\SI{6370000}{m}, the model sphere (\cref{sec:dominio}). With this declaration the coordinates that
eccodes reconstructs from the grid definition coincide with those of the model to within
\SI{2.009}{m} over the whole domain. Without it, the grid reconstructed by eccodes is displaced by
up to about \SI{1.1}{km} at the northern edge.

\section{Vertical levels}
\label{sec:livelli}

The 13 three-dimensional variables are on the same 13 pressure levels,
\SIlist{1000;925;850;700;600;500;400;300;250;200;150;100;50}{\hPa}. These are the levels
requested by the LAM AI project, a subset of the 37 pressure levels of ERA5 \parencite{era5doc}.
Of the others, the 8 above \SI{50}{\hPa} lie above the model top and cannot be produced.

The 246 messages are distributed over the level types of \cref{tab:tipi-livello}.

\begin{table}[htbp]
  \centering
  \small
  \caption{Level types of the 246 messages of an hourly file.}
  \label{tab:tipi-livello}
  \begin{tabularx}{\textwidth}{@{}lrX@{}}
    \toprule
    Level type & Messages & Content \\
    \midrule
    \campo{isobaricInhPa}      & 169 & 13 variables on 13 pressure levels \\
    \campo{heightAboveGround}  & 11  & wind at 10, 100 and \SI{200}{m}; temperature, dew-point temperature and relative humidity at \SI{2}{m}; \campo{10fg}; \campo{vwsh} \\
    \campo{surface}            & 47  & static fields, soil, snow, precipitation, runoff and surface radiation \\
    \campo{entireAtmosphere}   & 10  & column integrals and cloud covers \\
    \campo{nominalTop}         & 5   & radiation at the top of the atmosphere \\
    \campo{mostUnstableParcel} & 2   & \campo{mucape}, \campo{mucin} \\
    \campo{meanSea}            & 1   & \campo{msl} \\
    \campo{entireLake}         & 1   & \campo{dl} \\
    \bottomrule
  \end{tabularx}
\end{table}

\section{Time convention}
\label{sec:tempo}

Instantaneous fields use product definition template 4.0, and their reference time
(\campo{dataDate}, \campo{dataTime}) coincides with the validity time. The 19 accumulated fields
are total precipitation (\campo{tp}), rain (\campo{tirf}), snowfall (\campo{sf}), the
three runoffs (\campo{sro}, \campo{ssro}, \campo{ro}) and the 13 radiative fluxes at the surface
and at the top of the atmosphere. They use template 4.8, because the ECMWF definitions of these
parameters require the key \campo{typeOfStatisticalProcessing}, which does not exist in template
4.0 \parencite{eccodes}. As in ERA5 \parencite{era5doc}, every accumulated field covers the hour
ending at the validity time $H$; in the archive the reference time is $H-1$ and the accumulation
interval (\campo{stepRange}) is 0--1. The value is the difference between two consecutive hours of the same
run, for example $\campo{tp}(H) = [\campo{RAINNC}(H) - \campo{RAINNC}(H-1)]/1000$, with
\campo{RAINNC} in \si{mm} and \campo{tp} in \si{m} \parencite{issue70}. The twentieth field with
template 4.8 is \campo{10fg}, the maximum of the \SI{10}{m} wind, with reference time $H-1$,
declared interval 0--1 and \campo{stepType=max} (\cref{cap:raffica}).

The hourly file at 00 UTC of day D comes entirely, accumulated fields included, from the thirtieth
hour of the run started at 18 UTC on D$-2$; the accumulated fields are the difference between its
thirtieth and twenty-ninth hours \parencite{issue71,issue72}. The accumulated fields of the hourly
file at 01 UTC of D are the difference between the seventh and sixth hours of the run started at
18 UTC on D$-1$. The 24 validity times from 01 UTC of D to 00 UTC of D$+1$ all come from that run,
and the sum of their accumulated fields is the run's total from its sixth to its thirtieth hour.
When the wrfout file of a run's thirtieth hour is missing at the source, every message of the
following hourly file at 00 UTC consists of missing points; as of 14 October 2026 this happens in no
hourly file of the archive.

Since instantaneous fields have reference time $H$ and fields with template 4.8 have $H-1$, an
hourly file contains two reference times, and an index of the hourly files must be built on the
validity time (\campo{validityDate}, \campo{validityTime}), never on the reference time.
Telling an accumulated field from an instantaneous one requires \campo{stepType} (or
\campo{typeOfStatisticalProcessing}) and \campo{stepRange}. The key
\campo{generatingProcessIdentifier} is 128, the eccodes default for GRIB2, in all 246 messages,
and does not separate accumulated fields from instantaneous ones \parencite{issue53}.

\section{The run switch}
\label{sec:cambio-run}

The \emph{run switch} is the step, once a day, between the 00 UTC validity time of D, thirtieth hour
of the run started at 18 UTC on D$-2$, and the 01 UTC validity time of D, seventh hour of the run
started at 18 UTC on D$-1$. The jump between the two validity times is the difference between the
two runs at 00 UTC of D, thirtieth hour of the first and sixth of the second, plus the change over
the hour 00--01 UTC in the second. \Cref{tab:cambio-run} compares it with
one hour of change within the same run, over the full domain, at a single run switch, from 23 UTC on
18 January to 01 UTC on 19 January 2024 \parencite{issue72}.

\begin{table}[htbp]
  \centering
  \small
  \caption{Root-mean-square difference between two consecutive validity times, over the full
  domain, 18 and 19 January 2024.}
  \label{tab:cambio-run}
  \begin{tabular}{@{}lSSSSS@{}}
    \toprule
    Validity times & {\campo{2t} (\si{K})} & {\campo{sp} (\si{\hPa})} & {\SI{10}{m} wind speed (\si{m/s})} & {\campo{t} at \SI{850}{\hPa} (\si{K})} & {\campo{tcwv} (\si{kg/m^2})} \\
    \midrule
    23--00 UTC, same run & 0.66 & 0.50 & 1.13 & 0.48 & 0.89 \\
    00--01 UTC, run switch & 1.79 & 1.62 & 1.62 & 0.94 & 1.54 \\
    \bottomrule
  \end{tabular}
\end{table}

On 19 January 2024 the root-mean-square difference of \campo{2t} at the run switch is \num{2.7}
times that of one hour within the same run, and the maximum absolute difference between the two
validity times is \SI{11.0}{K}, against \SI{9.1}{K} between 23 and 00 UTC. A tendency or a time derivative computed across the interval from 00 to 01 UTC
contains this jump.
